\emph{Conventions.} As in Theorem \claimref{synthesis.field-strength-cancels-from-the-leading-null-singularity-of-hadamard-dirac-states-on-null-hypersurfaces}: units $\hbar = c = 1$, signature $(+,-,-,-)$, the physical null frame $\{k, n, f_{1}, f_{2}\}$ parallel along the generators, $k\cdot n = 1$, the canonical projector $\Lambda = \tfrac12 nk$, the area density $a$, the expansion $\theta = a'/a$, the transverse Hessian $M(u, u')$ of $-\sigma$ in the labels, the Van Vleck determinant $\Delta$, the canonical generator field $\psi_{\mathcal{H}} = \sqrt{2a}\,\Lambda\psi$, $K = L^{2}(\R^{2}, d^{2}y)\otimes V_{\Lambda}$ and $W_{0}(u, u') = \tfrac{i}{2\pi}(u - u' + i0)^{-1}$. The field algebra is the Dirac CAR algebra of Fact \claimref{synthesis.imported-dimock-dirac-car-algebra-on-a-gravitational-field}(a), whose anticommutator pairing is the conserved Cauchy inner product of the solutions (its (b)). A gauge-invariant quasi-free Hadamard state $\omega$ is one satisfying the microlocal condition of Fact \claimref{synthesis.imported-microlocal-hadamard-condition}(b); near any point its two-point functions are the Dirac operator applied to a Hadamard parametrix of the squared operator plus a smooth kernel (Fact \claimref{synthesis.imported-dirac-hadamard-form-propagates-and-adiabatic-states-are-hadamard}(a)), with leading coefficient $U = \Delta^{1/2}R$ (Fact \claimref{synthesis.imported-hadamard-parametrix-of-the-squared-dirac-operator-leads-with-van-vleck-root-times-spin-transport}(a), (b)).

\emph{Step 1 ((a), the bulk scaling limit).} This is Theorem \claimref{synthesis.hadamard-dirac-states-scale-to-the-flat-massless-vacuum-at-every-point-of-a-gravitational-field}(iii): in any gauge with $\hb(x_{0}) = 1$, for every $m_{0} \ge 0$ and whatever the Weyl and Ricci field strength at $x_{0}$, $l^{3}\omega(\psi(x_{0} + ly)\psi^{\dagger}(x_{0} + ly')) \to i\nabla D_{+}(y - y')e_{0}$ and the other ordering tends to $-i\nabla D_{-}(y - y')e_{0}$, and since $\omega$ is gauge-invariant and quasi-free all its $n$-point functions converge to those of the flat massless vacuum $\omega_{\mathrm{fl}}$.

\emph{Step 2 (the geometry of the horizon).} Let $\mathcal{H}$ be a local Rindler horizon through $x_{0}$ (Premise \claimref{synthesis.c-loc-local-horizon-premise} names these surfaces; nothing it asserts is used), or any piece of a null hypersurface through $x_{0}$ generated by affinely parametrised null geodesics. Shrink it to a piece in a geodesically and causally convex $N$ on which (S) holds: distinct generators are spacelike separated. (S) holds on every small enough piece (the transverse Hessian is positive near the diagonal) and on every local Rindler horizon, the null boundary of the past of a small two-surface element being achronal (Theorem \claimref{synthesis.restricted-dirac-hadamard-parametrix-on-a-null-hypersurface-is-the-null-vacuum-plus-a-bounded-kernel}, hypotheses; standard causal theory, cited). The piece lies in a smooth local foliation $\{\mathcal{H}_{v}\}$ by null hypersurfaces of the same kind (the null hypersurfaces orthogonal to a family of parallel two-surfaces; standard, cited from memory). On it, Theorem \claimref{synthesis.field-strength-cancels-from-the-leading-null-singularity-of-hadamard-dirac-states-on-null-hypersurfaces}(i) gives $\sigma(X(u, y' + \eta), X(u', y')) = -\tfrac12 M(\eta, \eta) + O(|\eta|^{3})$ and
\[
\frac{\Delta^{1/2}}{\sqrt{\det M}} = \bigl(a(u)\,a(u')\bigr)^{-1/2},\qquad \Bigl(\partial_{u} + \tfrac12\theta(u)\Bigr)\bigl(a(u)a(u')\bigr)^{-1/2} = 0,
\]
the second by $\theta = a'/a$ (the half-density law). This one identity is what makes both the step coefficient of Step 3 and the logarithmic coefficient of Step 5 vanish.

\emph{Step 3 (the generator anticommutation relations).} By Theorem \claimref{synthesis.dirac-states-restrict-to-null-hypersurfaces-through-the-flux-contraction-and-are-thin-layer-limits-under-the-generator-car}(i) the Dirac current is future causal, so the trace $\iota_{v}^{*}\Phi = \sqrt{2a_{v}}\,\Lambda\Phi|_{\mathcal{H}_{v}}$ on every leaf is a contraction, its adjoint defines the generator field inside the field algebra, $\psi_{\mathcal{H}}(F) = \psi(\iota_{\mathcal{H}}F)$, and every state restricts to bounded positive two-point operators with $W^{+}_{\mathcal{H},\omega} + W^{-}_{\mathcal{H},\omega} = \iota_{\mathcal{H}}^{*}\iota_{\mathcal{H}}$. By Theorem \claimref{synthesis.dirac-causal-propagator-between-nearby-null-leaves-tends-to-the-generator-car-on-every-geodesic-null-hypersurface}(i) the causal propagator restricted between $\mathcal{H}$ and $\mathcal{H}_{v}$ tends to $\delta(u - u')\delta^{2}(y - y')\Lambda$ as $v \to 0$ from either side, its step coefficient being $a(u')(\partial_{u} + \tfrac12\theta)(\Delta^{1/2}/\sqrt{\det M})\Lambda = 0$ by Step 2, for every mass. With the contraction this gives that theorem's (ii): $\iota_{\mathcal{H}}^{*}\iota_{\mathcal{H}} = 1$, the generator anticommutation relations $\{\psi_{\mathcal{H}}(u, y), \psi_{\mathcal{H}}(u', y')^{*}\} = \delta(u - u')\delta^{2}(y - y')\Lambda$, on every leaf. So (Iso) holds on the foliation, and $W^{-}_{\mathcal{H},\omega} = 1 - W^{+}_{\mathcal{H},\omega}$ for every state.

\emph{Step 4 ((b), the restriction of the states).} By Theorem \claimref{synthesis.dirac-states-restrict-to-null-hypersurfaces-through-the-flux-contraction-and-are-thin-layer-limits-under-the-generator-car}(iii-a), with no further hypothesis, the restricted two-point functions are iterated limits over thin layers of the bulk ones, and those of two gauge-invariant quasi-free Hadamard states differ by the pullback $D_{0,0}$ of the smooth difference of their bulk kernels; so $W^{\mathrm{par}}_{\mathcal{H}} := W^{+}_{\mathcal{H},\omega} - D^{\omega}_{0,0}$, with $D^{\omega}$ the smooth part of $\omega$'s kernel beyond the parametrix, is the same for every such state. By Theorem \claimref{synthesis.dirac-causal-propagator-between-nearby-null-leaves-tends-to-the-generator-car-on-every-geodesic-null-hypersurface}(iii)(a) the same Morse data give the thin-layer limit of the parametrix: $W_{0}\otimes\id_{K} + r$, with the same $i0$ from both sides of the leaf, no $\log|v|$ term (its coefficient is the jet combination of Step 2, zero), and, for each pair of transverse profiles, $r(u, u')$ a bounded function, bounded uniformly for profiles bounded in $C^{4}$ with support in a fixed compact set; by its (iii)(b),
\[
W^{+}_{\mathcal{H},\omega} = W_{0}\otimes\id_{K} + r + D^{\omega}_{0,0}.
\]
For $u_{0}$ interior and test functions $\varphi$, $\chi$, $\lambda\int\!\!\int\bar\varphi(x)W_{0}(u_{0} + \lambda x, u_{0} + \lambda x')\chi(x')\,dx\,dx' = \int\!\!\int\bar\varphi W_{0}\chi$ ($W_{0}$ is homogeneous of degree $-1$), while for a bounded or smooth kernel $b$, $|\lambda^{\varpi}\int\!\!\int\bar\varphi(x)\,b(u_{0} + \lambda x, u_{0} + \lambda x')\,\chi(x')\,dx\,dx'| \le \lambda^{\varpi}\sup|b|\,\|\varphi\|_{1}\|\chi\|_{1} \to 0$ for every $\varpi > 0$. So $\lambda^{\varpi}(W^{+}_{\mathcal{H},\omega} - W_{0}\otimes\id_{K})(u_{0} + \lambda\cdot, u_{0} + \lambda\cdot) \to 0$ for every $\varpi > 0$, uniformly in the profiles as stated, and $\lambda W^{+}_{\mathcal{H},\omega}(u_{0} + \lambda\cdot, u_{0} + \lambda\cdot) \to W_{0}\otimes\id_{K}$; by Step 3 the other ordering tends to $(\delta - W_{0})\otimes\id_{K}$. The restricted state is quasi-free and gauge-invariant, so its generator scaling limit at $x_{0}$ is the null vacuum $\omega_{0}$. This is (b).

\emph{Step 5 ((c), the $\epsilon$-regularised restriction).} Theorem \claimref{synthesis.restricted-dirac-hadamard-parametrix-on-a-null-hypersurface-is-the-null-vacuum-plus-a-bounded-kernel} treats the restriction in which versions 1 to 3 posed the premise, the $\epsilon \to 0$ limit of the transverse-smeared, $i0$-regularised parametrix: its numerator vanishes to second order in the label offset, the coefficient of $\log\epsilon$ is $2\pi(a a')^{1/2}(\partial_{u} + \tfrac12\theta)(\Delta^{1/2}/\sqrt{\det M})\Lambda$, zero by Step 2, the coefficient of $(u - u' + i0)^{-1}\delta^{2}(y - y')$ is that of $W_{0}\Lambda$ (Theorem \claimref{synthesis.field-strength-cancels-from-the-leading-null-singularity-of-hadamard-dirac-states-on-null-hypersurfaces}(ii)), and the rest is bounded: the restriction is $W_{0}\otimes\id_{K}$ plus a bounded kernel, of generator scaling degree at most $0$ (its (i) to (iii)). It depends on the time function $t$ by the smooth kernel $\tfrac{i}{2\pi}\partial_{u}\log\varphi_{t}\,\Lambda\,\delta^{2}(y - y')$, $t(X(u, y)) - t(X(u', y)) = (u - u')\varphi_{t}$ (its (ii)): on the flat null plane with $t = u + cu^{2}$ along the generators, $t_{u}/(t(u) - t(u')) - (u - u')^{-1} = c/(1 + c(u + u'))$ (checked exactly). The thin-layer limit of Step 4 does not depend on the time function (it is a limit of genuine pullbacks between distinct leaves). By Theorem \claimref{synthesis.dirac-causal-propagator-between-nearby-null-leaves-tends-to-the-generator-car-on-every-geodesic-null-hypersurface}(iii)(b) the two restrictions of the parametrix therefore differ by a bounded kernel containing that smooth term, nonzero when $dt(k)$ varies along the generators: they agree modulo a bounded, time-function-dependent kernel, not identically. In particular the clause of version 3, $\lambda^{\varpi}R^{\mathrm{par}}(u_{0} + \lambda\cdot, u_{0} + \lambda\cdot) \to 0$ with some $\varpi < 1$, holds for either restriction, with every $\varpi \in (0, 1)$.

\emph{Step 6 ((d), scaling and restriction commute; the local form of (N$_{+}$)).} By Theorem \claimref{synthesis.bulk-hadamard-states-restrict-to-the-generator-condition-on-flat-null-planes}(ii) at $m_{0} = 0$, the restriction of the bulk limit $\omega_{\mathrm{fl}}$ of Step 1 to the null plane through $x_{0}$ tangent to $\mathcal{H}$ is $W_{0}\otimes\id_{K}$ in the canonical normalisation: the generator scaling limit of Step 4 is the restriction of the bulk scaling limit, so scaling at $x_{0}$ and restriction to the generators commute, which is the clause (b) of version 2. The generator scaling limit being the null vacuum at $x_{0}$ is the local form of (N$_{+}$) at $x_{0}$ (in the sense of Premise \claimref{synthesis.n-plus-null-ground-state}, local form, and Theorem \claimref{synthesis.local-normality-at-the-edge-gives-the-null-ground-state}(ii); neither is used as a premise). Nothing in Steps 1 to 6 uses the field equations, vanishing expansion or shear, or the value of the field strength at $x_{0}$, and every step holds for every $m_{0} \ge 0$.

\emph{What remains cited.} The chain rests on the imported Facts named above (within them, the parts their notes mark as cited from memory: the Lichnerowicz identification of the connection of the squared Dirac operator, the four-dimensional leading coefficient and the bundle case of the fundamental solutions, Friedlander and B\"ar, Ginoux and Pf\"affle, Ch.~2), and, inside the proofs of the results it assembles, on standard inputs that are not imported: the existence of a local null foliation containing a small piece of a geodesic null hypersurface; the Morse lemma with parameters (H\"ormander, Vol.~III, App.~C, lemma number not checked) and the pullback of distributions (H\"ormander, Vol.~I, Thm.~8.2.4); the achronality of local Rindler horizons, which gives (S) there; the divergence theorem on a compact region with Lipschitz boundary in a globally hyperbolic field. The jump-relation route of the generator-CAR result is argued and not used. On light cones near the vertex the generator anticommutation relations of Step 3 agree with the characteristic Cauchy problem of Häfner and Nicolas (Fact \claimref{synthesis.imported-haefner-nicolas-dirac-characteristic-problem-on-light-cones-near-the-vertex}, supporting provenance, not used).

\emph{Exact checks (computed).} (1) The canonical projector: $\Lambda^{2} = \Lambda$, $\Lambda k = 0$, $k\Lambda = k$, $\Lambda pk = 2(p\cdot k)\Lambda$ and $(k\cdot e_{0})\Lambda pe_{0}\Lambda = (p\cdot k)\Lambda$ for four null directions and four regulators, and the null flux of the Dirac current $2(k\cdot e_{0})|\Lambda\psi|^{2}$ on two spinors. (2) The identity of Step 2 on the future light cone of a point of the flat field ($\Delta = 1$, expansion $2/u$), with $M$ computed from the Synge function by exact truncated series in a rational stereographic chart: $\sigma$ has no constant or linear term in the offset, $(\det M)^{2} = \det G(u)\det G(u')$ ($G$ the induced metric of the cut, $\det G = a^{2}$) at three point pairs, and $a(u) = u^{2}a(1)$, so the coefficient of Step 2 vanishes. (3) The time-function kernel of Step 5 on the flat null plane, bounded, scaling away at rate $\lambda$, and the dilation invariance of $W_{0}$. The analytic content (contraction, Morse reduction, limits, bounds) is in the assembled proofs and their checks.

\emph{Status.} Proved, given the imported Facts and the standard inputs listed. Transport-absent: free field theory on a Gauge Theory Gravity field; Gauge Theory Gravity is convenient (one spinor trivialisation, the projector algebra in the spacetime algebra).
